\chapter*{Abstract}
\addcontentsline{toc}{chapter}{Abstract}

Sepsis kills about eleven million people a year, nearly one death in five
worldwide, and in intensive care every hour of delayed treatment raises the
risk of dying. More than ninety machine-learning systems have been built to
raise the alarm earlier. The recent systematic evidence across them finds that
their apparent accuracy survives the move to a new hospital while their
clinical usefulness does not. This thesis asks how much of that reported
performance belongs to the model at all, and how much belongs to four choices
the analyst makes before any model is fitted: how sepsis is defined in the
record (the \emph{label}, which decides which patients count as septic and from
which hour), how the data are split, when the prediction is timed, and which
metric is reported. On \pcNStays{} adult ICU stays from MIMIC-IV, with the
models applied without retraining to a second database, three pre-registered
readings of the Sepsis-3 definition were crossed with six model
specifications, two split designs and two metrics, the key estimates carrying
stay-level bootstrap intervals and each declared family of comparisons
carrying a correction for its size. The premise did
not survive. Across the six specifications compared, the model choice moved
discrimination by \pcSpreadModelClass{} while the label moved it by
\pcSpreadLabel, roughly a tenth as much. But the two narrowest readings of
the same definition agreed on which patients were septic only at
$\kappa = \pcKappaVsBA$, so the label was deciding not how well a cohort could
be ranked but which cohort existed to be ranked at all. Its second consequence
is timing: because the standard onset rule fires at a documented clinical
action, the window before treatment contains no sepsis events by construction
rather than by observation, and re-timing the identical stays by the field's
alternative rule places \pcPreTreatOnsetsVarF{} events in it, while a label
built from physiology alone places \pcPreTreatOnsetsVarD. Every model's
discrimination is of the onset hour itself rather than of a future horizon, and
no evaluated operating point comes near the alert-burden benchmark this field
sets: the best costs \pcUtilityOptBurdenGbtB{} false alerts for every true one
against a benchmark of 1.4. Improving discrimination alone therefore does not
resolve the membership, timing and alert-burden limitations of a
treatment-anchored sepsis label. Clinical actions constitute its membership,
and its onset timing decides whether a prediction task defined before those
actions can be posed at all.

\vspace{0.4em}
\noindent\textbf{Keywords:} sepsis; MIMIC-IV; discrete-time survival analysis;
label definition; variance decomposition; clinical utility; TRIPOD+AI.
